\pdfoutput=1
%% Author: PGL  Porta Mana
%% Created: 2015-05-01T20:53:34+0200
%% Last-Updated: 2022-02-03T13:41:53+0100
%%%%%%%%%%%%%%%%%%%%%%%%%%%%%%%%%%%%%%%%%%%%%%%%%%%%%%%%%%%%%%%%%%%%%%%%%%%%
\newif\ifarxiv
\arxivfalse
\ifarxiv\pdfmapfile{+classico.map}\fi
\newif\ifafour
\afourfalse% true = A4, false = A5
\newif\iftypodisclaim % typographical disclaim on the side
\typodisclaimtrue
\newcommand*{\memfontfamily}{zplx}
\newcommand*{\memfontpack}{newpxtext}
\documentclass[\ifafour a4paper,12pt,\else a5paper,10pt,\fi%extrafontsizes,%
onecolumn,oneside,article,%french,italian,german,swedish,latin,
british%
]{memoir}
\newcommand*{\firstdraft}{17 December 2021}
\newcommand*{\firstpublished}{\firstdraft}
\newcommand*{\updated}{\ifarxiv***\else\today\fi}
\newcommand*{\propertitle}{Report on results and technical details of AD study\\{\large for Ingrid, Alexandra, \amp\ co.}%
}% title uses LARGE; set Large for smaller
\newcommand*{\pdftitle}{\propertitle}
\newcommand*{\headtitle}{AD report}
\newcommand*{\pdfauthor}{P.G.L.  Porta Mana}
\newcommand*{\headauthor}{Luca}
\newcommand*{\reporthead}{\iftrue\else Open Science Framework \href{https://doi.org/10.31219/osf.io/***}{\textsc{doi}:10.31219/osf.io/***}\fi}% Report number

%%%%%%%%%%%%%%%%%%%%%%%%%%%%%%%%%%%%%%%%%%%%%%%%%%%%%%%%%%%%%%%%%%%%%%%%%%%%
%%% Calls to packages (uncomment as needed)
%%%%%%%%%%%%%%%%%%%%%%%%%%%%%%%%%%%%%%%%%%%%%%%%%%%%%%%%%%%%%%%%%%%%%%%%%%%%

%\usepackage{pifont}

%\usepackage{fontawesome}

\usepackage[T1]{fontenc} 
\input{glyphtounicode} \pdfgentounicode=1

\usepackage[utf8]{inputenx}

%\usepackage{newunicodechar}
% \newunicodechar{Ĕ}{\u{E}}
% \newunicodechar{ĕ}{\u{e}}
% \newunicodechar{Ĭ}{\u{I}}
% \newunicodechar{ĭ}{\u{\i}}
% \newunicodechar{Ŏ}{\u{O}}
% \newunicodechar{ŏ}{\u{o}}
% \newunicodechar{Ŭ}{\u{U}}
% \newunicodechar{ŭ}{\u{u}}
% \newunicodechar{Ā}{\=A}
% \newunicodechar{ā}{\=a}
% \newunicodechar{Ē}{\=E}
% \newunicodechar{ē}{\=e}
% \newunicodechar{Ī}{\=I}
% \newunicodechar{ī}{\={\i}}
% \newunicodechar{Ō}{\=O}
% \newunicodechar{ō}{\=o}
% \newunicodechar{Ū}{\=U}
% \newunicodechar{ū}{\=u}
% \newunicodechar{Ȳ}{\=Y}
% \newunicodechar{ȳ}{\=y}

\newcommand*{\bmmax}{0} % reduce number of bold fonts, before font packages
\newcommand*{\hmmax}{0} % reduce number of heavy fonts, before font packages

\usepackage{textcomp}

%\usepackage[normalem]{ulem}% package for underlining
% \makeatletter
% \def\ssout{\bgroup \ULdepth=-.35ex%\UL@setULdepth
%  \markoverwith{\lower\ULdepth\hbox
%    {\kern-.03em\vbox{\hrule width.2em\kern1.2\p@\hrule}\kern-.03em}}%
%  \ULon}
% \makeatother

\usepackage{amsmath}

\usepackage{mathtools}
%\addtolength{\jot}{\jot} % increase spacing in multiline formulae
\setlength{\multlinegap}{0pt}

%\usepackage{empheq}% automatically calls amsmath and mathtools
%\newcommand*{\widefbox}[1]{\fbox{\hspace{1em}#1\hspace{1em}}}

%%%% empheq above seems more versatile than these:
%\usepackage{fancybox}
%\usepackage{framed}

% \usepackage[misc]{ifsym} % for dice
% \newcommand*{\diceone}{{\scriptsize\Cube{1}}}

\usepackage{amssymb}

\usepackage{amsxtra}

\usepackage[main=british]{babel}\selectlanguage{british}
%\newcommand*{\langnohyph}{\foreignlanguage{nohyphenation}}
\newcommand{\langnohyph}[1]{\begin{hyphenrules}{nohyphenation}#1\end{hyphenrules}}

\usepackage[autostyle=false,autopunct=false,english=british]{csquotes}
\setquotestyle{british}
\newcommand*{\defquote}[1]{`\,#1\,'}

% \makeatletter
% \renewenvironment{quotation}%
%                {\list{}{\listparindent 1.5em%
%                         \itemindent    \listparindent
%                         \rightmargin=1em   \leftmargin=1em
%                         \parsep        \z@ \@plus\p@}%
%                 \item[]\footnotesize}%
%                 {\endlist}
% \makeatother                


\usepackage{amsthm}
%% from https://tex.stackexchange.com/a/404680/97039
\makeatletter
\def\@endtheorem{\endtrivlist}
\makeatother

\newcommand*{\QED}{\textsc{q.e.d.}}
\renewcommand*{\qedsymbol}{\QED}
\theoremstyle{remark}
\newtheorem{note}{Note}
\newtheorem*{remark}{Note}
\newtheoremstyle{innote}{\parsep}{\parsep}{\footnotesize}{}{}{}{0pt}{}
\theoremstyle{innote}
\newtheorem*{innote}{}

\usepackage[shortlabels,inline]{enumitem}
\SetEnumitemKey{para}{itemindent=\parindent,leftmargin=0pt,listparindent=\parindent,parsep=0pt,itemsep=\topsep}
% \begin{asparaenum} = \begin{enumerate}[para]
% \begin{inparaenum} = \begin{enumerate*}
\setlist{itemsep=0pt,topsep=\parsep}
\setlist[enumerate,2]{label=\alph*.}
\setlist[enumerate]{label=\arabic*.,leftmargin=1.5\parindent}
\setlist[itemize]{leftmargin=1.5\parindent}
\setlist[description]{leftmargin=1.5\parindent}
% old alternative:
% \setlist[enumerate,2]{label=\alph*.}
% \setlist[enumerate]{leftmargin=\parindent}
% \setlist[itemize]{leftmargin=\parindent}
% \setlist[description]{leftmargin=\parindent}

\usepackage[babel,theoremfont,largesc]{newpxtext}

\usepackage[bigdelims,nosymbolsc%,smallerops % probably arXiv doesn't have it
]{newpxmath}
%\useosf
%\linespread{1.083}%
%\linespread{1.05}% widely used
\linespread{1.1}% best for text with maths
%% smaller operators for old version of newpxmath
\makeatletter
\def\re@DeclareMathSymbol#1#2#3#4{%
    \let#1=\undefined
    \DeclareMathSymbol{#1}{#2}{#3}{#4}}
%\re@DeclareMathSymbol{\bigsqcupop}{\mathop}{largesymbols}{"46}
%\re@DeclareMathSymbol{\bigodotop}{\mathop}{largesymbols}{"4A}
\re@DeclareMathSymbol{\bigoplusop}{\mathop}{largesymbols}{"4C}
\re@DeclareMathSymbol{\bigotimesop}{\mathop}{largesymbols}{"4E}
\re@DeclareMathSymbol{\sumop}{\mathop}{largesymbols}{"50}
\re@DeclareMathSymbol{\prodop}{\mathop}{largesymbols}{"51}
\re@DeclareMathSymbol{\bigcupop}{\mathop}{largesymbols}{"53}
\re@DeclareMathSymbol{\bigcapop}{\mathop}{largesymbols}{"54}
%\re@DeclareMathSymbol{\biguplusop}{\mathop}{largesymbols}{"55}
\re@DeclareMathSymbol{\bigwedgeop}{\mathop}{largesymbols}{"56}
\re@DeclareMathSymbol{\bigveeop}{\mathop}{largesymbols}{"57}
%\re@DeclareMathSymbol{\bigcupdotop}{\mathop}{largesymbols}{"DF}
%\re@DeclareMathSymbol{\bigcapplusop}{\mathop}{largesymbolsPXA}{"00}
%\re@DeclareMathSymbol{\bigsqcupplusop}{\mathop}{largesymbolsPXA}{"02}
%\re@DeclareMathSymbol{\bigsqcapplusop}{\mathop}{largesymbolsPXA}{"04}
%\re@DeclareMathSymbol{\bigsqcapop}{\mathop}{largesymbolsPXA}{"06}
\re@DeclareMathSymbol{\bigtimesop}{\mathop}{largesymbolsPXA}{"10}
%\re@DeclareMathSymbol{\coprodop}{\mathop}{largesymbols}{"60}
%\re@DeclareMathSymbol{\varprod}{\mathop}{largesymbolsPXA}{16}
\makeatother
%%
%% With euler font cursive for Greek letters - the [1] means 100% scaling
\DeclareFontFamily{U}{egreek}{\skewchar\font'177}%
\DeclareFontShape{U}{egreek}{m}{n}{<-6>s*[1]eurm5 <6-8>s*[1]eurm7 <8->s*[1]eurm10}{}%
\DeclareFontShape{U}{egreek}{m}{it}{<->s*[1]eurmo10}{}%
\DeclareFontShape{U}{egreek}{b}{n}{<-6>s*[1]eurb5 <6-8>s*[1]eurb7 <8->s*[1]eurb10}{}%
\DeclareFontShape{U}{egreek}{b}{it}{<->s*[1]eurbo10}{}%
\DeclareSymbolFont{egreeki}{U}{egreek}{m}{it}%
\SetSymbolFont{egreeki}{bold}{U}{egreek}{b}{it}% from the amsfonts package
\DeclareSymbolFont{egreekr}{U}{egreek}{m}{n}%
\SetSymbolFont{egreekr}{bold}{U}{egreek}{b}{n}% from the amsfonts package
% Take also \sum, \prod, \coprod symbols from Euler fonts
\DeclareFontFamily{U}{egreekx}{\skewchar\font'177}
\DeclareFontShape{U}{egreekx}{m}{n}{%
       <-7.5>s*[0.9]euex7%
    <7.5-8.5>s*[0.9]euex8%
    <8.5-9.5>s*[0.9]euex9%
    <9.5->s*[0.9]euex10%
}{}
\DeclareSymbolFont{egreekx}{U}{egreekx}{m}{n}
\DeclareMathSymbol{\sumop}{\mathop}{egreekx}{"50}
\DeclareMathSymbol{\prodop}{\mathop}{egreekx}{"51}
\DeclareMathSymbol{\coprodop}{\mathop}{egreekx}{"60}
\makeatletter
\def\sum{\DOTSI\sumop\slimits@}
\def\prod{\DOTSI\prodop\slimits@}
\def\coprod{\DOTSI\coprodop\slimits@}
\makeatother
% Greek letters not usually given in LaTeX.
\DeclareMathSymbol{\varpartial}{\mathalpha}{egreeki}{"40}
\DeclareMathSymbol{\partialup}{\mathalpha}{egreekr}{"40}
\DeclareMathSymbol{\alpha}{\mathalpha}{egreeki}{"0B}
\DeclareMathSymbol{\beta}{\mathalpha}{egreeki}{"0C}
\DeclareMathSymbol{\gamma}{\mathalpha}{egreeki}{"0D}
\DeclareMathSymbol{\delta}{\mathalpha}{egreeki}{"0E}
\DeclareMathSymbol{\epsilon}{\mathalpha}{egreeki}{"0F}
\DeclareMathSymbol{\zeta}{\mathalpha}{egreeki}{"10}
\DeclareMathSymbol{\eta}{\mathalpha}{egreeki}{"11}
\DeclareMathSymbol{\theta}{\mathalpha}{egreeki}{"12}
\DeclareMathSymbol{\iota}{\mathalpha}{egreeki}{"13}
\DeclareMathSymbol{\kappa}{\mathalpha}{egreeki}{"14}
\DeclareMathSymbol{\lambda}{\mathalpha}{egreeki}{"15}
\DeclareMathSymbol{\mu}{\mathalpha}{egreeki}{"16}
\DeclareMathSymbol{\nu}{\mathalpha}{egreeki}{"17}
\DeclareMathSymbol{\xi}{\mathalpha}{egreeki}{"18}
\DeclareMathSymbol{\omicron}{\mathalpha}{egreeki}{"6F}
\DeclareMathSymbol{\pi}{\mathalpha}{egreeki}{"19}
\DeclareMathSymbol{\rho}{\mathalpha}{egreeki}{"1A}
\DeclareMathSymbol{\sigma}{\mathalpha}{egreeki}{"1B}
 \DeclareMathSymbol{\tau}{\mathalpha}{egreeki}{"1C}
\DeclareMathSymbol{\upsilon}{\mathalpha}{egreeki}{"1D}
\DeclareMathSymbol{\phi}{\mathalpha}{egreeki}{"1E}
\DeclareMathSymbol{\chi}{\mathalpha}{egreeki}{"1F}
\DeclareMathSymbol{\psi}{\mathalpha}{egreeki}{"20}
\DeclareMathSymbol{\omega}{\mathalpha}{egreeki}{"21}
\DeclareMathSymbol{\varepsilon}{\mathalpha}{egreeki}{"22}
\DeclareMathSymbol{\vartheta}{\mathalpha}{egreeki}{"23}
\DeclareMathSymbol{\varpi}{\mathalpha}{egreeki}{"24}
\let\varrho\rho 
\let\varsigma\sigma
 \let\varkappa\kappa
\DeclareMathSymbol{\varphi}{\mathalpha}{egreeki}{"27}
%
\DeclareMathSymbol{\varAlpha}{\mathalpha}{egreeki}{"41}
\DeclareMathSymbol{\varBeta}{\mathalpha}{egreeki}{"42}
\DeclareMathSymbol{\varGamma}{\mathalpha}{egreeki}{"00}
\DeclareMathSymbol{\varDelta}{\mathalpha}{egreeki}{"01}
\DeclareMathSymbol{\varEpsilon}{\mathalpha}{egreeki}{"45}
\DeclareMathSymbol{\varZeta}{\mathalpha}{egreeki}{"5A}
\DeclareMathSymbol{\varEta}{\mathalpha}{egreeki}{"48}
\DeclareMathSymbol{\varTheta}{\mathalpha}{egreeki}{"02}
 \DeclareMathSymbol{\varIota}{\mathalpha}{egreeki}{"49}
\DeclareMathSymbol{\varKappa}{\mathalpha}{egreeki}{"4B}
\DeclareMathSymbol{\varLambda}{\mathalpha}{egreeki}{"03}
\DeclareMathSymbol{\varMu}{\mathalpha}{egreeki}{"4D}
\DeclareMathSymbol{\varNu}{\mathalpha}{egreeki}{"4E}
\DeclareMathSymbol{\varXi}{\mathalpha}{egreeki}{"04}
\DeclareMathSymbol{\varOmicron}{\mathalpha}{egreeki}{"4F}
\DeclareMathSymbol{\varPi}{\mathalpha}{egreeki}{"05}
\DeclareMathSymbol{\varRho}{\mathalpha}{egreeki}{"50}
\DeclareMathSymbol{\varSigma}{\mathalpha}{egreeki}{"06}
\DeclareMathSymbol{\varTau}{\mathalpha}{egreeki}{"54}
\DeclareMathSymbol{\varUpsilon}{\mathalpha}{egreeki}{"07}
\DeclareMathSymbol{\varPhi}{\mathalpha}{egreeki}{"08}
\DeclareMathSymbol{\varChi}{\mathalpha}{egreeki}{"58}
\DeclareMathSymbol{\varPsi}{\mathalpha}{egreeki}{"09}
\DeclareMathSymbol{\varOmega}{\mathalpha}{egreeki}{"0A} 
%
\DeclareMathSymbol{\Alpha}{\mathalpha}{egreekr}{"41}
\DeclareMathSymbol{\Beta}{\mathalpha}{egreekr}{"42}
\DeclareMathSymbol{\Gamma}{\mathalpha}{egreekr}{"00}
\DeclareMathSymbol{\Delta}{\mathalpha}{egreekr}{"01}
\DeclareMathSymbol{\Epsilon}{\mathalpha}{egreekr}{"45}
\DeclareMathSymbol{\Zeta}{\mathalpha}{egreekr}{"5A}
\DeclareMathSymbol{\Eta}{\mathalpha}{egreekr}{"48}
\DeclareMathSymbol{\Theta}{\mathalpha}{egreekr}{"02}
\DeclareMathSymbol{\Iota}{\mathalpha}{egreekr}{"49}
\DeclareMathSymbol{\Kappa}{\mathalpha}{egreekr}{"4B}
\DeclareMathSymbol{\Lambda}{\mathalpha}{egreekr}{"03}
\DeclareMathSymbol{\Mu}{\mathalpha}{egreekr}{"4D}
\DeclareMathSymbol{\Nu}{\mathalpha}{egreekr}{"4E}
\DeclareMathSymbol{\Xi}{\mathalpha}{egreekr}{"04}
\DeclareMathSymbol{\Omicron}{\mathalpha}{egreekr}{"4F}
\DeclareMathSymbol{\Pi}{\mathalpha}{egreekr}{"05}
\DeclareMathSymbol{\Rho}{\mathalpha}{egreekr}{"50}
\DeclareMathSymbol{\Sigma}{\mathalpha}{egreekr}{"06}
\DeclareMathSymbol{\Tau}{\mathalpha}{egreekr}{"54}
\DeclareMathSymbol{\Upsilon}{\mathalpha}{egreekr}{"07}
\DeclareMathSymbol{\Phi}{\mathalpha}{egreekr}{"08}
\DeclareMathSymbol{\Chi}{\mathalpha}{egreekr}{"58}
\DeclareMathSymbol{\Psi}{\mathalpha}{egreekr}{"09}
\DeclareMathSymbol{\Omega}{\mathalpha}{egreekr}{"0A}
%
\DeclareMathSymbol{\alphaup}{\mathalpha}{egreekr}{"0B}
\DeclareMathSymbol{\betaup}{\mathalpha}{egreekr}{"0C}
\DeclareMathSymbol{\gammaup}{\mathalpha}{egreekr}{"0D}
 \DeclareMathSymbol{\deltaup}{\mathalpha}{egreekr}{"0E}
\DeclareMathSymbol{\epsilonup}{\mathalpha}{egreekr}{"0F}
\DeclareMathSymbol{\zetaup}{\mathalpha}{egreekr}{"10}
\DeclareMathSymbol{\etaup}{\mathalpha}{egreekr}{"11}
\DeclareMathSymbol{\thetaup}{\mathalpha}{egreekr}{"12}
\DeclareMathSymbol{\iotaup}{\mathalpha}{egreekr}{"13}
\DeclareMathSymbol{\kappaup}{\mathalpha}{egreekr}{"14}
\DeclareMathSymbol{\lambdaup}{\mathalpha}{egreekr}{"15}
\DeclareMathSymbol{\muup}{\mathalpha}{egreekr}{"16}
\DeclareMathSymbol{\nuup}{\mathalpha}{egreekr}{"17}
\DeclareMathSymbol{\xiup}{\mathalpha}{egreekr}{"18}
\DeclareMathSymbol{\omicronup}{\mathalpha}{egreekr}{"6F}
  \DeclareMathSymbol{\piup}{\mathalpha}{egreekr}{"19}
\DeclareMathSymbol{\rhoup}{\mathalpha}{egreekr}{"1A}
\DeclareMathSymbol{\sigmaup}{\mathalpha}{egreekr}{"1B}
\DeclareMathSymbol{\tauup}{\mathalpha}{egreekr}{"1C}
\DeclareMathSymbol{\upsilonup}{\mathalpha}{egreekr}{"1D}
\DeclareMathSymbol{\phiup}{\mathalpha}{egreekr}{"1E}
\DeclareMathSymbol{\chiup}{\mathalpha}{egreekr}{"1F}
\DeclareMathSymbol{\psiup}{\mathalpha}{egreekr}{"20}
\DeclareMathSymbol{\omegaup}{\mathalpha}{egreekr}{"21}
\DeclareMathSymbol{\varepsilonup}{\mathalpha}{egreekr}{"22}
\DeclareMathSymbol{\varthetaup}{\mathalpha}{egreekr}{"23}
\DeclareMathSymbol{\varpiup}{\mathalpha}{egreekr}{"24}
\let\varrhoup\rhoup 
\let\varsigmaup\sigmaup
\let\varkappaup\kappaup
\DeclareMathSymbol{\varphiup}{\mathalpha}{egreekr}{"27}
% Greek letters not usually given in LaTeX.

%\usepackage%[scaled=0.9]%
%{classico}%  Optima as sans-serif font
\usepackage{classico} % \renewcommand\sfdefault{URWClassico-TLF}
\DeclareMathAlphabet{\mathsf}  {T1}{\sfdefault}{m}{sl}
\SetMathAlphabet{\mathsf}{bold}{T1}{\sfdefault}{b}{sl}
%\newcommand*{\mathte}[1]{\textbf{\textit{\textsf{#1}}}}
% Upright sans-serif math alphabet
% \DeclareMathAlphabet{\mathsu}  {T1}{\sfdefault}{m}{n}
% \SetMathAlphabet{\mathsu}{bold}{T1}{\sfdefault}{b}{n}

% DejaVu Mono as typewriter text
\usepackage[scaled=0.84]{DejaVuSansMono}

\usepackage{mathdots}

\usepackage[usenames]{xcolor}
% Tol (2012) colour-blind-, print-, screen-friendly colours, alternative scheme; Munsell terminology
\definecolor{mypurpleblue}{RGB}{68,119,170}
\definecolor{myblue}{RGB}{102,204,238}
\definecolor{mygreen}{RGB}{34,136,51}
\definecolor{myyellow}{RGB}{204,187,68}
\definecolor{myred}{RGB}{238,102,119}
\definecolor{myredpurple}{RGB}{170,51,119}
\definecolor{mygrey}{RGB}{187,187,187}
% Tol (2012) colour-blind-, print-, screen-friendly colours; Munsell terminology
% \definecolor{lbpurple}{RGB}{51,34,136}
% \definecolor{lblue}{RGB}{136,204,238}
% \definecolor{lbgreen}{RGB}{68,170,153}
% \definecolor{lgreen}{RGB}{17,119,51}
% \definecolor{lgyellow}{RGB}{153,153,51}
% \definecolor{lyellow}{RGB}{221,204,119}
% \definecolor{lred}{RGB}{204,102,119}
% \definecolor{lpred}{RGB}{136,34,85}
% \definecolor{lrpurple}{RGB}{170,68,153}
\definecolor{lgrey}{RGB}{221,221,221}
%\newcommand*\mycolourbox[1]{%
%\colorbox{mygrey}{\hspace{1em}#1\hspace{1em}}}
\colorlet{shadecolor}{lgrey}

\usepackage{bm}

\usepackage{microtype}

\usepackage[backend=biber,mcite,%subentry,
citestyle=authoryear-comp,bibstyle=pglpm-authoryear,autopunct=false,sorting=ny,sortcites=false,natbib=false,maxcitenames=2,maxbibnames=8,minbibnames=8,giveninits=true,uniquename=false,uniquelist=false,maxalphanames=1,block=space,hyperref=true,defernumbers=false,useprefix=true,sortupper=false,language=british,parentracker=false]{biblatex}
\DeclareSortingTemplate{ny}{\sort{\field{sortname}\field{author}\field{editor}}\sort{\field{year}}}
\iffalse\makeatletter%%% replace parenthesis with brackets
\newrobustcmd*{\parentexttrack}[1]{%
  \begingroup
  \blx@blxinit
  \blx@setsfcodes
  \blx@bibopenparen#1\blx@bibcloseparen
  \endgroup}
\AtEveryCite{%
  \let\parentext=\parentexttrack%
  \let\bibopenparen=\bibopenbracket%
  \let\bibcloseparen=\bibclosebracket}
\makeatother\fi
\DefineBibliographyExtras{british}{\def\finalandcomma{\addcomma}}
\renewcommand*{\finalnamedelim}{\addspace\amp\space}
% \renewcommand*{\finalnamedelim}{\addcomma\space}
\renewcommand*{\textcitedelim}{\addcomma\space}
% \setcounter{biburlnumpenalty}{1} % to allow url breaks anywhere
% \setcounter{biburlucpenalty}{0}
% \setcounter{biburllcpenalty}{1}
\DeclareDelimFormat{multicitedelim}{\addsemicolon\addspace\space}
\DeclareDelimFormat{compcitedelim}{\addsemicolon\addspace\space}
\DeclareDelimFormat{postnotedelim}{\addspace}
\ifarxiv\else\addbibresource{portamanabib.bib}\fi
\renewcommand{\bibfont}{\footnotesize}
%\appto{\citesetup}{\footnotesize}% smaller font for citations
\defbibheading{bibliography}[\bibname]{\section*{#1}\addcontentsline{toc}{section}{#1}%\markboth{#1}{#1}
}
\newcommand*{\citep}{\footcites}
\newcommand*{\citey}{\footcites}%{\parencites*}
\newcommand*{\ibid}{\unspace\addtocounter{footnote}{-1}\footnotemark{}}
%\renewcommand*{\cite}{\parencite}
%\renewcommand*{\cites}{\parencites}
\providecommand{\href}[2]{#2}
\providecommand{\eprint}[2]{\texttt{\href{#1}{#2}}}
\newcommand*{\amp}{\&}
% \newcommand*{\citein}[2][]{\textnormal{\textcite[#1]{#2}}%\addtocategory{extras}{#2}
% }
\newcommand*{\citein}[2][]{\textnormal{\textcite[#1]{#2}}%\addtocategory{extras}{#2}
}
\newcommand*{\citebi}[2][]{\textcite[#1]{#2}%\addtocategory{extras}{#2}
}
\newcommand*{\subtitleproc}[1]{}
\newcommand*{\chapb}{ch.}
%
%\def\UrlOrds{\do\*\do\-\do\~\do\'\do\"\do\-}%
\def\myUrlOrds{\do\0\do\1\do\2\do\3\do\4\do\5\do\6\do\7\do\8\do\9\do\a\do\b\do\c\do\d\do\e\do\f\do\g\do\h\do\i\do\j\do\k\do\l\do\m\do\n\do\o\do\p\do\q\do\r\do\s\do\t\do\u\do\v\do\w\do\x\do\y\do\z\do\A\do\B\do\C\do\D\do\E\do\F\do\G\do\H\do\I\do\J\do\K\do\L\do\M\do\N\do\O\do\P\do\Q\do\R\do\S\do\T\do\U\do\V\do\W\do\X\do\Y\do\Z}%
\makeatletter
%\g@addto@macro\UrlSpecials{\do={\newline}}
\g@addto@macro{\UrlBreaks}{\myUrlOrds}
\makeatother
\newcommand*{\arxiveprint}[1]{%
\href{https://arxiv.org/abs/#1}{arXiv:\allowbreak\nolinkurl{#1}}%
}
\newcommand*{\mparceprint}[1]{%
\href{http://www.ma.utexas.edu/mp_arc-bin/mpa?yn=#1}{mp\_arc:\allowbreak\nolinkurl{#1}}%
}
\newcommand*{\haleprint}[1]{%
\href{https://hal.archives-ouvertes.fr/#1}{\textsc{hal}:\allowbreak\nolinkurl{#1}}%
}
\newcommand*{\philscieprint}[1]{%
\href{http://philsci-archive.pitt.edu/archive/#1}{PhilSci:\allowbreak\nolinkurl{#1}}%
}
\newcommand*{\doi}[1]{%
\href{https://doi.org/#1}{\textsc{doi}:\allowbreak\nolinkurl{#1}}%
}
\newcommand*{\biorxiveprint}[1]{%
bioRxiv \doi{10.1101/#1}%
}
\newcommand*{\osfeprint}[1]{%
Open Science Framework \doi{10.31219/osf.io/#1}%
}

\usepackage{graphicx}

%\usepackage{wrapfig}

%\usepackage{tikz-cd}

\PassOptionsToPackage{hyphens}{url}\usepackage[hypertexnames=false,pdfencoding=unicode,psdextra]{hyperref}

\usepackage[depth=4]{bookmark}
\hypersetup{colorlinks=true,bookmarksnumbered,pdfborder={0 0 0.25},citebordercolor={0.2667 0.4667 0.6667},citecolor=mypurpleblue,linkbordercolor={0.6667 0.2 0.4667},linkcolor=myredpurple,urlbordercolor={0.1333 0.5333 0.2},urlcolor=mygreen,breaklinks=true,pdftitle={\pdftitle},pdfauthor={\pdfauthor}}
% \usepackage[vertfit=local]{breakurl}% only for arXiv
\providecommand*{\urlalt}{\href}

\usepackage[british]{datetime2}
\DTMnewdatestyle{mydate}%
{% definitions
\renewcommand*{\DTMdisplaydate}[4]{%
\number##3\ \DTMenglishmonthname{##2} ##1}%
\renewcommand*{\DTMDisplaydate}{\DTMdisplaydate}%
}
\DTMsetdatestyle{mydate}

%%%%%%%%%%%%%%%%%%%%%%%%%%%%%%%%%%%%%%%%%%%%%%%%%%%%%%%%%%%%%%%%%%%%%%%%%%%%
%%% Layout. I do not know on which kind of paper the reader will print the
%%% paper on (A4? letter? one-sided? double-sided?). So I choose A5, which
%%% provides a good layout for reading on screen and save paper if printed
%%% two pages per sheet. Average length line is 66 characters and page
%%% numbers are centred.
%%%%%%%%%%%%%%%%%%%%%%%%%%%%%%%%%%%%%%%%%%%%%%%%%%%%%%%%%%%%%%%%%%%%%%%%%%%%
\ifafour\setstocksize{297mm}{210mm}%{*}% A4
\else\setstocksize{210mm}{5.5in}%{*}% 210x139.7
\fi
\settrimmedsize{\stockheight}{\stockwidth}{*}
\setlxvchars[\normalfont] %313.3632pt for a 66-characters line
\setxlvchars[\normalfont]
% \setlength{\trimtop}{0pt}
% \setlength{\trimedge}{\stockwidth}
% \addtolength{\trimedge}{-\paperwidth}
%\settrims{0pt}{0pt}
% The length of the normalsize alphabet is 133.05988pt - 10 pt = 26.1408pc
% The length of the normalsize alphabet is 159.6719pt - 12pt = 30.3586pc
% Bringhurst gives 32pc as boundary optimal with 69 ch per line
% The length of the normalsize alphabet is 191.60612pt - 14pt = 35.8634pc
\ifafour\settypeblocksize{*}{32pc}{1.618} % A4
%\setulmargins{*}{*}{1.667}%gives 5/3 margins % 2 or 1.667
\else\settypeblocksize{*}{26pc}{1.618}% nearer to a 66-line newpx and preserves GR
\fi
\setulmargins{*}{*}{1}%gives equal margins
\setlrmargins{*}{*}{*}
\setheadfoot{\onelineskip}{2.5\onelineskip}
\setheaderspaces{*}{2\onelineskip}{*}
\setmarginnotes{2ex}{10mm}{0pt}
\checkandfixthelayout[nearest]
%%% End layout
%% this fixes missing white spaces
%\pdfmapline{+dummy-space <dummy-space.pfb}
%\pdfinterwordspaceon% seems to add a white margin to Sumatrapdf

%%% Sectioning
\newcommand*{\asudedication}[1]{%
{\par\centering\textit{#1}\par}}
\newenvironment{acknowledgements}{\section*{Thanks}\addcontentsline{toc}{section}{Thanks}}{\par}
\makeatletter\renewcommand{\appendix}{\par
  \bigskip{\centering
   \interlinepenalty \@M
   \normalfont
   \printchaptertitle{\sffamily\appendixpagename}\par}
  \setcounter{section}{0}%
  \gdef\@chapapp{\appendixname}%
  \gdef\thesection{\@Alph\c@section}%
  \anappendixtrue}\makeatother
\counterwithout{section}{chapter}
\setsecnumformat{\upshape\csname the#1\endcsname\quad}
\setsecheadstyle{\large\bfseries\sffamily%
\centering}
\setsubsecheadstyle{\bfseries\sffamily%
\raggedright}
%\setbeforesecskip{-1.5ex plus 1ex minus .2ex}% plus 1ex minus .2ex}
%\setaftersecskip{1.3ex plus .2ex }% plus 1ex minus .2ex}
%\setsubsubsecheadstyle{\bfseries\sffamily\slshape\raggedright}
%\setbeforesubsecskip{1.25ex plus 1ex minus .2ex }% plus 1ex minus .2ex}
%\setaftersubsecskip{-1em}%{-0.5ex plus .2ex}% plus 1ex minus .2ex}
\setsubsecindent{0pt}%0ex plus 1ex minus .2ex}
\setparaheadstyle{\bfseries\sffamily%
\raggedright}
\setcounter{secnumdepth}{2}
\setlength{\headwidth}{\textwidth}
\newcommand{\addchap}[1]{\chapter*[#1]{#1}\addcontentsline{toc}{chapter}{#1}}
\newcommand{\addsec}[1]{\section*{#1}\addcontentsline{toc}{section}{#1}}
\newcommand{\addsubsec}[1]{\subsection*{#1}\addcontentsline{toc}{subsection}{#1}}
\newcommand{\addpara}[1]{\paragraph*{#1.}\addcontentsline{toc}{subsubsection}{#1}}
\newcommand{\addparap}[1]{\paragraph*{#1}\addcontentsline{toc}{subsubsection}{#1}}

%%% Headers, footers, pagestyle
\copypagestyle{manaart}{plain}
\makeheadrule{manaart}{\headwidth}{0.5\normalrulethickness}
\makeoddhead{manaart}{%
{\footnotesize%\sffamily%
\scshape\headauthor}}{}{{\footnotesize\sffamily%
\headtitle}}
\makeoddfoot{manaart}{}{\thepage}{}
\newcommand*\autanet{\includegraphics[height=\heightof{M}]{autanet.pdf}}
\definecolor{mygray}{gray}{0.333}
\iftypodisclaim%
\ifafour\newcommand\addprintnote{\begin{picture}(0,0)%
\put(245,149){\makebox(0,0){\rotatebox{90}{\tiny\color{mygray}\textsf{This
            document is designed for screen reading and
            two-up printing on A4 or Letter paper}}}}%
\end{picture}}% A4
\else\newcommand\addprintnote{\begin{picture}(0,0)%
\put(176,112){\makebox(0,0){\rotatebox{90}{\tiny\color{mygray}\textsf{This
            document is designed for screen reading and
            two-up printing on A4 or Letter paper}}}}%
\end{picture}}\fi%afourtrue
\makeoddfoot{plain}{}{\makebox[0pt]{\thepage}\addprintnote}{}
\else
\makeoddfoot{plain}{}{\makebox[0pt]{\thepage}}{}
\fi%typodisclaimtrue
\makeoddhead{plain}{\scriptsize\reporthead}{}{}
% \copypagestyle{manainitial}{plain}
% \makeheadrule{manainitial}{\headwidth}{0.5\normalrulethickness}
% \makeoddhead{manainitial}{%
% \footnotesize\sffamily%
% \scshape\headauthor}{}{\footnotesize\sffamily%
% \headtitle}
% \makeoddfoot{manaart}{}{\thepage}{}

\pagestyle{manaart}

\setlength{\droptitle}{-3.9\onelineskip}
\pretitle{\begin{center}\LARGE\sffamily%
\bfseries}
\posttitle{\bigskip\end{center}}

\makeatletter\newcommand*{\atf}{\includegraphics[totalheight=\heightof{@}]{atblack.png}}\makeatother
\providecommand{\affiliation}[1]{\textsl{\textsf{\footnotesize #1}}}
\providecommand{\epost}[1]{\texttt{\footnotesize\textless#1\textgreater}}
\providecommand{\email}[2]{\href{mailto:#1ZZ@#2 ((remove ZZ))}{#1\protect\atf#2}}
%\providecommand{\email}[2]{\href{mailto:#1@#2}{#1@#2}}

\preauthor{\vspace{-0.5\baselineskip}\begin{center}
\normalsize\sffamily%
\lineskip  0.5em}
\postauthor{\par\end{center}}
\predate{\DTMsetdatestyle{mydate}\begin{center}\footnotesize}
\postdate{\end{center}\vspace{-\medskipamount}}

\setfloatadjustment{figure}{\footnotesize}
\captiondelim{\quad}
\captionnamefont{\footnotesize\sffamily%
}
\captiontitlefont{\footnotesize}
%\firmlists*
\midsloppy
% handling orphan/widow lines, memman.pdf
% \clubpenalty=10000
% \widowpenalty=10000
% \raggedbottom
% Downes, memman.pdf
\clubpenalty=9996
\widowpenalty=9999
\brokenpenalty=4991
\predisplaypenalty=10000
\postdisplaypenalty=1549
\displaywidowpenalty=1602
\raggedbottom

\paragraphfootnotes
\setlength{\footmarkwidth}{2ex}
% \threecolumnfootnotes
%\setlength{\footmarksep}{0em}
\footmarkstyle{\textsuperscript{%\color{myred}
\scriptsize\bfseries#1}~}
%\footmarkstyle{\textsuperscript{\color{myred}\scriptsize\bfseries#1}~}
%\footmarkstyle{\textsuperscript{[#1]}~}

\selectlanguage{british}\frenchspacing

%%%%%%%%%%%%%%%%%%%%%%%%%%%%%%%%%%%%%%%%%%%%%%%%%%%%%%%%%%%%%%%%%%%%%%%%%%%%
%%% Paper's details
%%%%%%%%%%%%%%%%%%%%%%%%%%%%%%%%%%%%%%%%%%%%%%%%%%%%%%%%%%%%%%%%%%%%%%%%%%%%
\title{\propertitle}
\author{%
\hspace*{\stretch{1}}%
%% uncomment if additional authors present
% \parbox{0.5\linewidth}%\makebox[0pt][c]%
% {\protect\centering ***\\%
% \footnotesize\epost{\email{***}{***}}}%
% \hspace*{\stretch{1}}%
\parbox{1\linewidth}%\makebox[0pt][c]%
{\protect\centering Luca  \href{https://orcid.org/0000-0002-6070-0784}{\protect\includegraphics[scale=0.16]{orcid_32x32.png}}%\\
\footnotesize%Western Norway University of Applied Sciences%
\quad\epost{\email{pgl}{portamana.org}}}%
% Mohn Medical Imaging and Visualization Centre, Dept of Computer science, Electrical Engineering and Mathematical Sciences, Western Norway University of Applied Sciences, Bergen, Norway
%% uncomment if additional authors present
% \hspace*{\stretch{1}}%
% \parbox{0.5\linewidth}%\makebox[0pt][c]%
% {\protect\centering ***\\%
% \footnotesize\epost{\email{***}{***}}}%
\hspace*{\stretch{1}}%
}

%\date{Draft of \today\ (first drafted \firstdraft)}
\date{\firstpublished; updated \updated}

%%%%%%%%%%%%%%%%%%%%%%%%%%%%%%%%%%%%%%%%%%%%%%%%%%%%%%%%%%%%%%%%%%%%%%%%%%%%
%%% Macros @@@
%%%%%%%%%%%%%%%%%%%%%%%%%%%%%%%%%%%%%%%%%%%%%%%%%%%%%%%%%%%%%%%%%%%%%%%%%%%%

% Common ones - uncomment as needed
%\providecommand{\nequiv}{\not\equiv}
%\providecommand{\coloneqq}{\mathrel{\mathop:}=}
%\providecommand{\eqqcolon}{=\mathrel{\mathop:}}
%\providecommand{\varprod}{\prod}
\newcommand*{\de}{\partialup}%partial diff
\newcommand*{\pu}{\piup}%constant pi
\newcommand*{\delt}{\deltaup}%Kronecker, Dirac
%\newcommand*{\eps}{\varepsilonup}%Levi-Civita, Heaviside
%\newcommand*{\riem}{\zetaup}%Riemann zeta
%\providecommand{\degree}{\textdegree}% degree
%\newcommand*{\celsius}{\textcelsius}% degree Celsius
%\newcommand*{\micro}{\textmu}% degree Celsius
\newcommand*{\I}{\mathrm{i}}%imaginary unit
\newcommand*{\e}{\mathrm{e}}%Neper
\newcommand*{\di}{\mathrm{d}}%differential
%\newcommand*{\Di}{\mathrm{D}}%capital differential
%\newcommand*{\planckc}{\hslash}
%\newcommand*{\avogn}{N_{\textrm{A}}}
%\newcommand*{\NN}{\bm{\mathrm{N}}}
%\newcommand*{\ZZ}{\bm{\mathrm{Z}}}
%\newcommand*{\QQ}{\bm{\mathrm{Q}}}
\newcommand*{\RR}{\bm{\mathrm{R}}}
%\newcommand*{\CC}{\bm{\mathrm{C}}}
%\newcommand*{\nabl}{\bm{\nabla}}%nabla
%\DeclareMathOperator{\lb}{lb}%base 2 log
%\DeclareMathOperator{\tr}{tr}%trace
%\DeclareMathOperator{\card}{card}%cardinality
%\DeclareMathOperator{\im}{Im}%im part
%\DeclareMathOperator{\re}{Re}%re part
%\DeclareMathOperator{\sgn}{sgn}%signum
%\DeclareMathOperator{\ent}{ent}%integer less or equal to
%\DeclareMathOperator{\Ord}{O}%same order as
%\DeclareMathOperator{\ord}{o}%lower order than
%\newcommand*{\incr}{\triangle}%finite increment
\newcommand*{\defd}{\coloneqq}
\newcommand*{\defs}{\eqqcolon}
%\newcommand*{\Land}{\bigwedge}
%\newcommand*{\Lor}{\bigvee}
%\newcommand*{\lland}{\DOTSB\;\land\;}
%\newcommand*{\llor}{\DOTSB\;\lor\;}
\newcommand*{\limplies}{\mathbin{\Rightarrow}}%implies
%\newcommand*{\suchthat}{\mid}%{\mathpunct{|}}%such that (eg in sets)
%\newcommand*{\with}{\colon}%with (list of indices)
%\newcommand*{\mul}{\times}%multiplication
%\newcommand*{\inn}{\cdot}%inner product
%\newcommand*{\dotv}{\mathord{\,\cdot\,}}%variable place
%\newcommand*{\comp}{\circ}%composition of functions
%\newcommand*{\con}{\mathbin{:}}%scal prod of tensors
%\newcommand*{\equi}{\sim}%equivalent to 
\renewcommand*{\asymp}{\simeq}%equivalent to 
%\newcommand*{\corr}{\mathrel{\hat{=}}}%corresponds to
%\providecommand{\varparallel}{\ensuremath{\mathbin{/\mkern-7mu/}}}%parallel (tentative symbol)
\renewcommand*{\le}{\leqslant}%less or equal
\renewcommand*{\ge}{\geqslant}%greater or equal
%\DeclarePairedDelimiter\clcl{[}{]}
%\DeclarePairedDelimiter\clop{[}{[}
%\DeclarePairedDelimiter\opcl{]}{]}
%\DeclarePairedDelimiter\opop{]}{[}
\DeclarePairedDelimiter\abs{\lvert}{\rvert}
%\DeclarePairedDelimiter\norm{\lVert}{\rVert}
\DeclarePairedDelimiter\set{\{}{\}} %}
%\DeclareMathOperator{\pr}{P}%probability
\newcommand*{\p}{\mathrm{p}}%probability
\renewcommand*{\P}{\mathrm{P}}%probability
%\newcommand*{\E}{\mathrm{E}}
%% The "\:" space is chosen to correctly separate inner binary and external rels
\renewcommand*{\|}[1][]{\nonscript\:#1\vert\nonscript\:\mathopen{}}
%\DeclarePairedDelimiterX{\cp}[2]{(}{)}{#1\nonscript\:\delimsize\vert\nonscript\:\mathopen{}#2}
%\DeclarePairedDelimiterX{\ct}[2]{[}{]}{#1\nonscript\;\delimsize\vert\nonscript\:\mathopen{}#2}
%\DeclarePairedDelimiterX{\cs}[2]{\{}{\}}{#1\nonscript\:\delimsize\vert\nonscript\:\mathopen{}#2}
%\newcommand*{\+}{\lor}
%\renewcommand{\*}{\land}
%% symbol = for equality statements within probabilities
%% from https://tex.stackexchange.com/a/484142/97039
% \newcommand*{\eq}{\mathrel{\!=\!}}
% \let\texteq\=
% \renewcommand*{\=}{\TextOrMath\texteq\eq}
% \newcommand*{\eq}[1][=]{\mathrel{\!#1\!}}
\newcommand*{\mo}[1][=]{\mathord{\,#1\,}}
%%
\newcommand*{\sect}{\S}% Sect.~
\newcommand*{\sects}{\S\S}% Sect.~
\newcommand*{\chap}{ch.}%
\newcommand*{\chaps}{chs}%
\newcommand*{\bref}{ref.}%
\newcommand*{\brefs}{refs}%
%\newcommand*{\fn}{fn}%
\newcommand*{\eqn}{eq.}%
\newcommand*{\eqns}{eqs}%
\newcommand*{\fig}{fig.}%
\newcommand*{\figs}{figs}%
\newcommand*{\vs}{{vs}}
\newcommand*{\eg}{{e.g.}}
\newcommand*{\etc}{{etc.}}
\newcommand*{\ie}{{i.e.}}
%\newcommand*{\ca}{{c.}}
\newcommand*{\foll}{{ff.}}
%\newcommand*{\viz}{{viz}}
\newcommand*{\cf}{{cf.}}
%\newcommand*{\Cf}{{Cf.}}
%\newcommand*{\vd}{{v.}}
\newcommand*{\etal}{{et al.}}
%\newcommand*{\etsim}{{et sim.}}
%\newcommand*{\ibid}{{ibid.}}
%\newcommand*{\sic}{{sic}}
%\newcommand*{\id}{\mathte{I}}%id matrix
%\newcommand*{\nbd}{\nobreakdash}%
%\newcommand*{\bd}{\hspace{0pt}}%
%\def\hy{-\penalty0\hskip0pt\relax}
%\newcommand*{\labelbis}[1]{\tag*{(\ref{#1})$_\text{r}$}}
%\newcommand*{\mathbox}[2][.8]{\parbox[t]{#1\columnwidth}{#2}}
%\newcommand*{\zerob}[1]{\makebox[0pt][l]{#1}}
\newcommand*{\tprod}{\mathop{\textstyle\prod}\nolimits}
\newcommand*{\tsum}{\mathop{\textstyle\sum}\nolimits}
%\newcommand*{\tint}{\begingroup\textstyle\int\endgroup\nolimits}
%\newcommand*{\tland}{\mathop{\textstyle\bigwedge}\nolimits}
%\newcommand*{\tlor}{\mathop{\textstyle\bigvee}\nolimits}
%\newcommand*{\sprod}{\mathop{\textstyle\prod}}
%\newcommand*{\ssum}{\mathop{\textstyle\sum}}
%\newcommand*{\sint}{\begingroup\textstyle\int\endgroup}
%\newcommand*{\sland}{\mathop{\textstyle\bigwedge}}
%\newcommand*{\slor}{\mathop{\textstyle\bigvee}}
%\newcommand*{\T}{^\transp}%transpose
%%\newcommand*{\QEM}%{\textnormal{$\Box$}}%{\ding{167}}
%\newcommand*{\qem}{\leavevmode\unskip\penalty9999 \hbox{}\nobreak\hfill
%\quad\hbox{\QEM}}

%%%%%%%%%%%%%%%%%%%%%%%%%%%%%%%%%%%%%%%%%%%%%%%%%%%%%%%%%%%%%%%%%%%%%%%%%%%%
%%% Custom macros for this file @@@
%%%%%%%%%%%%%%%%%%%%%%%%%%%%%%%%%%%%%%%%%%%%%%%%%%%%%%%%%%%%%%%%%%%%%%%%%%%%
\definecolor{notecolour}{RGB}{68,170,153}
%\newcommand*{\puzzle}{\maltese}
\newcommand*{\puzzle}{{\fontencoding{U}\fontfamily{fontawesometwo}\selectfont\symbol{225}}}
\newcommand*{\wrench}{{\fontencoding{U}\fontfamily{fontawesomethree}\selectfont\symbol{114}}}
\newcommand*{\pencil}{{\fontencoding{U}\fontfamily{fontawesometwo}\selectfont\symbol{210}}}
\newcommand{\mynote}[1]{ {\color{notecolour}\puzzle\ #1}}

\newcommand*{\widebar}[1]{{\mkern1.5mu\skew{2}\overline{\mkern-1.5mu#1\mkern-1.5mu}\mkern 1.5mu}}

% \newcommand{\explanation}[4][t]{%\setlength{\tabcolsep}{-1ex}
% %\smash{
% \begin{tabular}[#1]{c}#2\\[0.5\jot]\rule{1pt}{#3}\\#4\end{tabular}}%}
% \newcommand*{\ptext}[1]{\text{\small #1}}
%\DeclareMathOperator*{\argsup}{arg\,sup}
\newcommand*{\dob}{degree of belief}
\newcommand*{\dobs}{degrees of belief}
\newcommand*{\bit}{\textrm{Sh}}
\newcommand*{\narrow}{\texttt{narrow}}
\newcommand*{\broad}{\texttt{broad}}
\newcommand*{\all}{\texttt{all}}
\newcommand*{\IX}{I_{\setminus X}}
%%% Custom macros end @@@

%%%%%%%%%%%%%%%%%%%%%%%%%%%%%%%%%%%%%%%%%%%%%%%%%%%%%%%%%%%%%%%%%%%%%%%%%%%%
%%% Beginning of document
%%%%%%%%%%%%%%%%%%%%%%%%%%%%%%%%%%%%%%%%%%%%%%%%%%%%%%%%%%%%%%%%%%%%%%%%%%%%
%\firmlists
\begin{document}
\captiondelim{\quad}\captionnamefont{\footnotesize}\captiontitlefont{\footnotesize}
\selectlanguage{british}\frenchspacing
\maketitle

%%%%%%%%%%%%%%%%%%%%%%%%%%%%%%%%%%%%%%%%%%%%%%%%%%%%%%%%%%%%%%%%%%%%%%%%%%%%
%%% Abstract
%%%%%%%%%%%%%%%%%%%%%%%%%%%%%%%%%%%%%%%%%%%%%%%%%%%%%%%%%%%%%%%%%%%%%%%%%%%%
\iffalse\abstractrunin
\abslabeldelim{}
\renewcommand*{\abstractname}{}
\setlength{\absleftindent}{0pt}
\setlength{\absrightindent}{0pt}
\setlength{\abstitleskip}{-\absparindent}
\begin{abstract}\labelsep 0pt%
  \noindent ***
% \\\noindent\emph{\footnotesize Note: Dear Reader
%     \amp\ Peer, this manuscript is being peer-reviewed by you. Thank you.}
% \par%\\[\jot]
% \noindent
% {\footnotesize PACS: ***}\qquad%
% {\footnotesize MSC: ***}%
%\qquad{\footnotesize Keywords: ***}
\end{abstract}
\fi\selectlanguage{british}\frenchspacing

%%%%%%%%%%%%%%%%%%%%%%%%%%%%%%%%%%%%%%%%%%%%%%%%%%%%%%%%%%%%%%%%%%%%%%%%%%%%
%%% Epigraph
%%%%%%%%%%%%%%%%%%%%%%%%%%%%%%%%%%%%%%%%%%%%%%%%%%%%%%%%%%%%%%%%%%%%%%%%%%%%
% \asudedication{\small ***}
% \vspace{\bigskipamount}
% \setlength{\epigraphwidth}{.7\columnwidth}
% %\epigraphposition{flushright}
% \epigraphtextposition{flushright}
% %\epigraphsourceposition{flushright}
% \epigraphfontsize{\footnotesize}
% \setlength{\epigraphrule}{0pt}
% %\setlength{\beforeepigraphskip}{0pt}
% %\setlength{\afterepigraphskip}{0pt}
% \epigraph{\emph{text}}{source}



%%%%%%%%%%%%%%%%%%%%%%%%%%%%%%%%%%%%%%%%%%%%%%%%%%%%%%%%%%%%%%%%%%%%%%%%%%%%
%%% BEGINNING OF MAIN TEXT
%%%%%%%%%%%%%%%%%%%%%%%%%%%%%%%%%%%%%%%%%%%%%%%%%%%%%%%%%%%%%%%%%%%%%%%%%%%%

Here are the main results about the \enquote{predictive power} of features.
The technical details of the inference and method are explained in
\sect~\ref{sec:techdet}.

Here I report the results for Alexandra's FAQ-based features. I'll report
the results for Ingrid's study soon (some calculations are still ongoing).

All material is available at
\url{https://github.com/pglpm/ADBayes}.

\section{Main results}
\label{sec:main_res}

Our general problem: to prognose the future onset of AD as opposed to
stable MCI -- predictand binary variate \texttt{Subgroup\_num\_} -- given
the set of twelve features: \texttt{AGE}, \texttt{RAVLT\_immediate},
\texttt{AVDEL30MIN\_neuro}, \texttt{AVDELTOT\_neuro},
\texttt{TRAASCOR\_neuro}, \texttt{TRABSCOR\_neuro},
\texttt{CATANIMSC\_neuro}, \texttt{GDTOTAL\_gds}, \texttt{LRHHC\_n\_long},
\texttt{LRLV\_n\_long}, \texttt{FAQ}. In the following I use shorter names
for them.

The \enquote{predictive power} of a set of features is measured using the
\emph{mutual information} between that set of features and the predictand
variate, measured in \emph{Shannons} (symbol \enquote{$\bit$}, see
\textcite{iso2008c}). The mutual information can in this case range from
$0\;\bit$, representing a complete lack of predictive power (features and
predictand are independent, so we can just as well flip a coin to make our
prognoses); to $1\;\bit$, representing perfect, deterministic prediction.

The operational meaning of mutual
information is quickly explained in \sect~\ref{sec:joint} and somewhat
more extensively and with references in \sect~\ref{sec:mutualionfo}.

The results are reported for three different calculation set-ups, denoted
\narrow, \broad, \all. The first two use only the training dataset and have
different prior smoothness preferences for the inference: \narrow\ and
\broad. The third uses \all\ datapoints, with a broad prior smoothness
preference. More details are given in
\sects~\ref{sec:understand}--\ref{sec:maths}. It should be noted that the
probability calculation should always use all data, and no division between
training and test set is necessary (probability theory is implicitly
already doing all such possible divisions).


\subsection{Results for individual features}
\label{sec:indiv_features}

This is the mutual information for each feature if it were used
\emph{individually} to make the prognosis, ranked
from highest to lowest:
\begin{center}
  \begin{tabular}[c]{l|ccc}
    feature &\multicolumn{3}{c}{mutual information/Sh} \\
    & \narrow & \broad & \all \\
\hline
    \texttt{AVDEL30MIN} & 0.13 & 0.1 & 0.09\\
    \texttt{RAVLT} & 0.12 & 0.09 & 0.08\\
    \texttt{FAQ} & 0.10 & 0.08 & 0.06\\
    \texttt{AVDELTOT} & 0.07 & 0.05 & 0.06\\
    \texttt{LRHHC} & 0.04 & 0.03 & 0.03\\
    \texttt{TRABSCOR} & 0.03 & 0.03 & 0.03\\
    \texttt{CATANIMSC} & 0.03 & 0.02 & 0.02\\
    \texttt{TRAASCOR} & 0.03 & 0.02 & 0.01\\
    \texttt{LRLV} & 0.01 & 0.01 & 0.01\\
    \texttt{AGE} & 0.00 & 0.00 & 0.00\\
    \texttt{GDTOTAL} & 0.00 & 0.00 & 0.00
  \end{tabular}                            
\end{center}
The rank is the same in all computations set-ups; the values agree within
the first significant digit. These values do not only give us a ranking,
but also an estimate of the predictive power. See the next section for what these
numbers actually mean.

\subsection{Results for joint feature set}
\label{sec:joint}

This is the mutual information for the set of features, used \emph{jointly}:
\begin{center}
  \begin{tabular}[c]{ccc}
    \narrow & \broad & \all \\
    0.31\;\bit & 0.20\;\bit & 0.18\;\bit
  \end{tabular}
\end{center}

What does a value such as \enquote{0.2\;\bit} actually mean?

A mutual information $I=0.2\;\bit$ between the predictand and a set of
features gives us a range of the number of correct prognoses we can
expect to do about the predictand by using that set of features.

In a set of $N$ new prognoses, we have an approximate lower bound of
$(1+I)\,N/2 \pm \sqrt{(1-I)\,N}$ correct ones (true positives~+ true
negatives) on average, and an upper bound of $C\,N \pm 2\sqrt{(1-C)\,C\,N}$
correct ones (the interval is for a 95\% certainty) on average, with $C$
given by the formula
\begin{equation}
  \label{eq:mutualinfo_est}
 C = \frac{1}{2}\Bigl[1+\sqrt{1-(1-I)^{4/3}}\Bigr] \ .
\end{equation}
So, with a mutual information of 0.2\;\bit, in 100 new prognoses we can
expect  that between $60\pm 9$ and $75\pm 9$
will be correct. Analogously, a mutual information of $0.3\;\bit$ means
from $65 \pm 8$ to $81\pm 8$ correct prognoses in 100 new cases.

Note that this value is an implicit characteristic of this set of features:
it is a mathematical result from information theory that no predictive
algorithm can do better than this by using these features as predictors.
Algorithms that do less are not fully using the information in the
predictors; algorithms that happen to do more have only had a stroke of
luck.

\subsection{\enquote{Importance} of individual features when used jointly}
\label{sec:importance_features}

We must define what we mean by \enquote{importance} of a feature in the
joint prediction. Information is a very non-additive quantity: we cannot
say \enquote*{this feature contributes $x$ to the information, that feature
  contributes $y$} and so on, with $x+y+\dotsb$ adding up to the total
information. Said otherwise, the \enquote{importance} or
\enquote{contribution} of a feature is \emph{context-dependent}.

For example, imagine that that we have three features that jointly give us
a given amount of predictive power. It may happen that if we do not use the
first feature, our predictions are just as good; its contribution seems to
be zero. If we do not use the second feature, our predictions are also just
as good: its contribution seems to be zero as well. And yet, if we do not
use neither the first nor the second feature, our predictions become worse.
So \enquote{$0+0\ne 0$} in a manner of speaking. A figurative explanation
of this phenomenon is as follow: the first feature provides pieces of
information $AB$, the second feature provides $BC$, the third provides
$AC$. Jointly they provide $ABC$. If we drop the first, we still have $ABC$
from the remaining two. If we drop the second, we still have $ABC$. But if
we drop the first and the second, we only have $AB$; our prediction become
worse because of lack of the piece of information $C$.

\bigskip

From these considerations, my approach to somehow quantifying the
\enquote{importance} of feature $X$ is as follow:
\begin{enumerate}
\item calculate the mutual info of all features, used jointly; denote the
  result by $I$
\item calculate the mutual info of all features but excluding $X$; denote
  the result by $\IX$
\item calculate the relative decrease of mutual info we suffer when we
  remove $X$ from the features:
  \begin{equation}
    \label{eq:rel_decrease_MI_X}
    \Delta\IX \defd  \frac{I - \IX}{I}
    \equiv 1-\frac{\IX}{I}
  \end{equation}
\end{enumerate}
Note that by the properties of mutual info we always have $I \ge
\IX$, so $0 \le \Delta\IX \le 1$.

I define the \enquote{importance} of a feature $X$ as the relative decrease
in mutual information we suffer when $X$ is removed from the whole set of
twelve features used for prediction. It seems a sensible intuitive definition to
me, with values that have a concrete operational meaning (see below).

With this definition, here are the \enquote{importances} of the individual features,
ranked from highest to lowest according to the \all\ setup:
\begin{center}
  \begin{tabular}[c]{l|ccc}
    feature &\multicolumn{3}{c}{$\Delta I_{\setminus\text{feature}}$} \\
    & \narrow & \broad & \all \\
\hline
\texttt{TRABSCOR} & 14\% & 12\% & 13\%\\
\texttt{FAQ} & 13\% & 10\% & 9\%\\
\texttt{TRAASCOR} & 7\% & 4\% & 4\%\\
\texttt{AVDEL30MIN} & 4\% & 4\% & 4\%\\
\texttt{RAVLT} & 3\% & 1\% & 3\%\\
\texttt{AVDELTOT} & 2\% & 2\% & 2\%\\
\texttt{LRHHC} & 1\% & 1\% & 1\%\\
\texttt{CATANIMSC} & 1\% & 1\% & 0.7\%\\
\texttt{GDTOTAL} & 0.7\% & 0.1\% & 0.5\%\\
\texttt{LRLV} & 0.0\% & 0.3\% & 0.3\%\\
\texttt{AGE} & 0.4\% & 0.5\% & 0.3\%
  \end{tabular}
\end{center}
Note that the percentages do not add up to 100\%, nor should they, owing to
the reasons given above.

We see that omitting either \texttt{TRABSCOR\_neuro} or \texttt{FAQ} from
the twelve features would reduce the mutual information by roughly 10\%;
that is, from a value of $0.2\;\bit$ to $0.18\;\bit$. This also means that
in 100 new prognoses we would drop from a best case of $75\pm 9$ correct
ones to a best case of around $74\pm 9$. In clinical-importance terms this
difference is not small: it means on average more than 10\,000 additional
incorrect predictions (more exactly around 12\,500) every million
prognoses.

\subsection{Further remarks}
\label{sec:remarks}

\textbullet\ The mutual informations estimated above are calculated
assuming that \emph{we will not receive further training data}. It is also
possible to calculate an estimate of what the mutual informations would be
if we had a very large number of training data. I'm currently setting up
the code to calculate this estimate.

\medskip

\noindent\textbullet\ It is possible to calculate a more precise forecast
of how many correct predictions are possible in new prognoses. Let me know
if this is of interest.

\medskip

\noindent\textbullet\ As shown in the enclosed plots (all coming from the \all\ setup), many
features have two clearly different population distributions for the future
AD or MCI patients. This means that something is already at work affecting
those features. However, the two population distributions have very large
overlaps, so the feature does not help very much in making predictions for
a single individual, as clear also from the mutual informations for the
individual features.

\medskip

\noindent\textbullet\ A further analysis of direct and inverse conditional
frequencies in the population suggest that the dataset might have some
biases for some features. These biases, however, could be overcome by
appropriately combining inverse and direct predictions from the features.
This point and topic would require a very extensive discussion, so I won't
write any further details here.

\medskip

\noindent\textbullet\ As discussed in \sect~\ref{sec:maths}, the calculation also allows us to
count the number of peaks in the joint distribution of predictand and
features, possibly hinting at the presence of subpopulations. A very rough
estimate gives around 10 main peaks. If this is of interest, a more exact
calculation can be done.

\medskip


\noindent\textbullet\ Regarding \enquote{explanation}, I categorically refuse to say that the
mutual info or some other score of a feature \enquote{explains} anything.
For me, statistical and probabilistic analyses and metrics do not
\enquote{explain}, they only \emph{fit} or correlate. Using
\enquote{explanation} in this context is an intentional or unintentional
way of making it sound as if we have discovered more than we actually have
(the deceiving use of the word \enquote{explanation} for what's just a
\emph{fit} has been a very bad habit of statisticians since the 1960s or
70s). Explanations can only be given by hardcore physics, chemistry,
biology; but different levels of fit can of course suggest, or be evidence
for, different explanations formulated within those disciplines.


\section{Technical details}
\label{sec:techdet}

\subsection{Technical terminology}
\label{sec:terminology}

In the literature, the probabilistic calculations made here are usually
called \emph{Bayesian nonparametric density regression}. This terminology
can be psychologically misleading because it sounds as if we are making
peculiar assumptions, when in reality we are not: this is the most general
probabilistic inference possible. Here's a simple explanation of the
terminology.

\emph{Density regression} stands opposed to \emph{functional regression}.
Functional regression makes the assumption that there exists a physical
functional dependence between predictand and features, contaminated by
(usually gaussian) noise. Examples of algorithms that do functional
regression under the hood: linear regression, generalized linear models,
support vector machines, gaussian processes, neural networks. Density
regression does not assume that a functional relationship exists (it
doesn't exclude it either). It simply evaluates the probabilities of
observing different values of the predictand, given values of the features;
that is, it evaluates  probability \emph{densities}, hence the name.
Examples of algorithms that do density regression: logistic regression,
random forests (I believe).

\emph{Nonparametric} stands opposed to \emph{parametric}. Parametric
regression makes the assumption that the function or density has a specific
form or shape (line, exponential; gaussian; and so on). Examples of
parametric algorithms: linear regression, generalized linear models,
support vector machines, logistic regression. Nonparametric regression does
not make any such assumption; it considers all possible forms and shapes.
Examples of nonparametric algorithms: gaussian processes, neural networks,
random forests (I believe).

\emph{Nonparametric density regression} is therefore the most general
assumption-free inference we can possibly make.

\subsection{Intuitive understanding}
\label{sec:understand}

In simplified terms, probability theory in this inference first considers
every possible frequency distribution of predictand given features, for
\emph{all future} patients. Such frequency distribution also tells us what
we can guess about our next patient, given his/her features.

For example, suppose that by some powerful technology we knew that, among
\emph{all} -- or a billion of -- future patients that have feature
$\texttt{FAQ}=8$, a percentage of 68\% of them will develop AD and 32\% of
them will stay MCI. If our next patient has $\texttt{FAQ}=8$, then he/she
must be either one of the 680 millions who will develop AD, or one of of
the 320 millions who will not. So there's a 68\% probability that he/she
will develop AD. This consideration is made not only for \texttt{FAQ} but
for all 12 features jointly, and all their possible values. This simplified
understanding allows us to quickly interpret, for example, the solid red
line in the plot of \texttt{FAQ} \vs\ \enquote{probability of AD} enclosed
at the end of the present report.

\smallskip

Of course we do not know what the true frequency distribution of predictand
and features will be. Probability theory considers each possible frequency
distribution and attaches a probability -- a weight, if you like -- to it.
This weight, called \enquote{posterior}, comes from two contributions:
\begin{itemize}[wide]
\item[\emph{\enquote{Likelihood}}]\!: How much the distribution fits the data
  we already have. This is simply the joint probability of our data,
  assuming that frequency distribution were the true one.
\item[\emph{\enquote{Prior}}]\!: An initial weight based on biomedical
  considerations of how the true frequency could most likely look like.
\end{itemize}
In the present case we give a greater initial weight to distributions that
are smooth, without lots of sudden jumps or a zigzag shape, although they
may have as many large and small peaks as they like. This is biologically
very plausible, owing to continuity of physical and chemical processes.
Note that we are not excluding discontinuities: we only consider them less
plausible at first.

The total probability or weight is closer to the prior one when we have
very few data; but as the number of data increases, the preferences
expressed in the prior become less and less important. Below it's explained
what kind of prior preferences correspond to the \narrow\ and \broad\
setups.

The first result of this inference is therefore a probability distribution
over all possible frequency distributions for future patients. In other
words, how their \enquote{true} population distribution look like.

The second result is a probability for the next new patient we observe:
given his/her features, the probability is an average over all candidate
frequency distributions, weighted by their posteriors. This is summarized
in the following (slightly informal) formula:
\begin{subequations}  \label{eq:definetti}
  \begin{gather}
    \p(Y \|X) =    \int f(Y \| X) \ \p(f \| \textrm{\small data})\ \di f
    \\
    \shortintertext{with}
    \p(f \| \textrm{\small data}) \propto
    f(\textrm{\small data})\ \p(f \| H)
  \end{gather}
\end{subequations}
where $f(Y\|X)$ is the frequency of predictand value $Y$ among patients
with features $X$; the probability of our data given such frequency is
$f(\textrm{\small data})$; and $\p(f \| H)$ is the prior weight of such
frequency.


In summary, the probability calculation is guessing what the true
distribution of predictand and features is, for the full population, and it
also gives our uncertainty about such guess. It uses this guess to make
predictions about the next and all subsequent patients.

In passing, it can be proven that the posterior mathematically includes an
average of \emph{all possible} $k$-fold cross-validations, for all $k$, and
all possible divisions of the data into \enquote{training} and
\enquote{test} sets \citep{portamana2019b}. This gives an idea of why
probability calculations are so enormously computationally expensive.

\subsection{Brief overview of mathematical aspects}
\label{sec:maths}

From a mathematical point of view the main problem is how to represent a
generic frequency distribution $f$ presented in the intuitive explanation
above. There are an infinity of mathematically equivalent representations.
Here we use the representation as an \emph{infinite mixture of product
  kernels}: a generic joint frequency distribution $f(Y, X, Z, \dotsc)$ is
represented as
\begin{equation}
  \label{eq:freq_repr}
  f(Y, X_{1}, X_{2}, \dotsc) =
  \sum_{k=1}^{\infty} w_{k}\ Q_{s_{k}}(Y)\ Q_{a_{k}}(X)\ Q_{b_{k}}(Z) \dotsm
\end{equation}
with an infinite set of coefficients
$\set{w_{k},s_{k},a_{k},b_{k},\dotsc}$. These coefficients effectively are
a coordinate system on the infinite-dimensional manifold of frequency
distributions \citep[esp.\ Part~VII]{choquetbruhatetal1977_r1996}. The
various $Q$ distributions are gaussian for continuous variates, binomial
for integer variates, and Bernoulli for binary variates.

This representation allows for quick calculations of conditional and
marginal frequencies. For more precise mathematical details see
\textcite{dunsonetal2011}, and \textcite{rasmussen1999} who only considers
gaussian kernels though.

This mathematical representation also allows, in principle, an
interpretation of the frequency distribution in terms of clusters, and an
interpretation of the population in terms of subpopulations. Such
interpretations, however, are arbitrary and contentious if we do not first
justify why the subpopulations should have distributions like the $Q$s and
not some other distributions -- say, why gaussian and not Cauchy or Gumbel.
I therefore do not trust any such subpopulation interpretation of the
formula above; its meaning is only in the final sum.

However, the number of clusters with high weights $w_{k}$ does give us a
very rough idea of the number of major and minor peaks of the frequency
distribution, and therefore a vague idea of the number of possible
subpopulations.

In the representation above, the prior probability over the frequencies is
represented by a probability density over the coefficients
$\set{w_{k},s_{k},a_{k},b_{k},\dotsc}$. This probability density determines
our initial guess about the smoothness of the true frequency
distribution.

The \narrow\ and \broad\ setup roughly correspond to two the following
guesses:
\begin{itemize}[wide]
\item[\narrow]\!: with 97\% probability, the widths of individual peaks can
  be two orders of magnitude broader or narrower than the average width
\item[\broad]\!: with 61\% probability, the widths of individual peaks can
  be two orders of magnitude broader or narrower than the average width
\end{itemize}
The general probability distribution for such peaks is a compound gamma
distribution. The $\narrow$ setup therefore favours a frequency
distribution that does not have too narrow or too broad peaks; the $\broad$
setup gives much more freedom. The results show that the data overwhelm
either choice, so our conclusions are robust to variation in this kind of
preference.

\bigskip

The computation of the probability for the different frequencies is done by
Markov-chain Monte Carlo sampling, with Gibbs sampling; see
\textcite{neal1993} for an exhaustive discussion, and again
\textcite{dunsonetal2011}, \textcite{rasmussen1999} and references therein
for details, and \textcite[\chap~29]{mackay1995_r2005} for an overview. The
computation has converged when the Markov chain has become stationary;
stationarity has been checked for various quantities derived from the
frequency distribution, such as the joint probability of the data, the
means and variances, joint covariances, and full-dimensional covariance.
Convergence required roughly 48--60 hours of computation.



\subsection{Entropy and mutual information}
\label{sec:mutualionfo}

For a comprehensive discussion of the following concepts see the first
chapters of \textcite{mackay1995_r2005}, especially \chap~4, and also his
lectures
% \footnote{\href{https://www.youtube.com/watch?v=BCiZc0n6COY\&list=PLruBu5BI5n4aFpG32iMbdWoRVAA-Vcso6}{\texttt{\doiurl{https://www.youtube.com/watch?v=BCiZc0n6COY&list=PLruBu5BI5n4aFpG32iMbdWoRVAA-Vcso6}}}}
\footnote{\url{https://www.youtube.com/watch?v=BCiZc0n6COY&list=PLruBu5BI5n4aFpG32iMbdWoRVAA-Vcso6}},
besides \textcite{shannon1948} and the very accessible summary in
\sect~14.7 of \textcite{pressetal1986_r2007}. Other references about its
predictive meaning and practical use are for example
\textcite{lindley1956,kullback1959_r1978,woodward1953_r1964}. These notions
are their units are also international standards: see \textcite{iso2008c},
items 13-23--13-42.

The \emph{entropy} of the distribution for a variate $Y$, here denoted
$H(Y)$, represents our uncertainty about future occurrences of that
variate. For a binary variate it assumes a value between $0\;\bit$
(complete certainty and exact prediction) and $1\;\bit$ (complete
uncertainty). A value of $H$ (in Shannons) means that we are on average
uncertain about  $H\,N$ future outcomes. This means that we will on average make
\begin{multline}
  \label{eq:entropy_est}
    \text{from}\quad
    (1-H/2)\,N \pm \sqrt{H\,N}\\
    \shoveleft{
    \text{to}\quad    C\,N \pm 2\sqrt{(1-C)\,C\,N}
      \quad\text{correct predictions,}
      }
    \\[1\jot]
    \text{\small with}\quad C = \frac{1}{2}\Bigl[1+\sqrt{1-H^{4/3}}\Bigr] \ .
\end{multline}
The entropy is also related to the size of the typical
set of different sequences of values we can see in $N$ occurrences: the
size is $2^{H}$.

When we know the \emph{specific} value of a feature (or features) $X$, we
can usually make better predictions about $Y$, reflected by a sharper
probability distribution for $Y$ given the specific value $x$. The entropy
of this distribution is denoted $H(Y\|x)$. The meaning is as before.

In considering the \emph{general} predictive power of the feature $X$, some
important considerations are necessary. For example, it can happen that our
predictions are almost perfect (entropy of 0\;\bit) when the value of $X$
is in some range $R_{\text{good}}$; whereas their are very bad (entropy of
1\;\bit) when the value is in the complementary range $R_{\text{bad}}$. Is
then the feature $X$ a good or a bad predictor? The answer heavily depends
\emph{on which values of $X$ we'll encounter in our future applications}.
If all future applications typically have values of $X$ in
$R_{\text{good}}$, then $X$ is a good predictor; if they have values in
$R_{\text{bad}}$, then $X$ is a bad predictor; for intermediate cases we
have intermediate predictive power. Clearly this depends on the
\emph{probability distribution of future occurrences of $X$}, $p(X)$.

The \emph{conditional entropy} of $Y$ given features $X$ is exactly the
weighted combination of the entropies given each value of $X$, weighted by
their frequency of occurrence:
\begin{equation}
  \label{eq:cond_entr}
  H(Y|X) \defd \sum_{x} H(Y\|x)\ p(X)
\end{equation}
Note that any kind of metric of predictive performance that neglects $p(X)$
is therefore missing a very important point; its results will be very
hit-or-miss. Conditional entropy correctly keeps this factor into account,
instead; this is the reason why its values can be given immediate
operational meaning in terms of prediction.


Finally, the \emph{mutual information}  between $Y$ and $X$ is simply the
difference between our uncertainty when we do not know $X$ and when we do:
\begin{equation}
  \label{eq:mutualinfo}
  I(Y,X) \defd H(Y) - H(Y\|X).
\end{equation}
From this definition and the operational meaning of the entropy we also
obtain a concrete operational meaning for the mutual information. In
particular, if we are completely uncertain about $Y$ in absence of
features, $H(Y)=1\;\bit$, then a mutual information $I$ means that the
features $X$ allow us to make on average 
\begin{multline}
  \label{eq:entropy_est}
    \text{from}\quad (1+I)\,N/2 \pm \sqrt{(1-I)\,N} \\
    \text{to}\quad   C\,N \pm 2\sqrt{(1-C)\,C\,N}
    \quad \text{correct predictions,}
    \\[1\jot]
    \text{\small with}\quad C = \frac{1}{2}\Bigl[1+\sqrt{1-(1-I)^{4/3}}\Bigr] \ .
\end{multline}
out of $N$ cases, as discussed in \sect~\ref{sec:joint}. The exact number
depends on the distribution of features (which can be calculated from the
probabilistic analysis).


% \[ \color{mypurpleblue}\bm{a} \color{myredpurple}\mathbin{\bm{\land}}\color{mypurpleblue} \bm{b} \]


%%%% examples use empheq
%   \begin{empheq}[left={\mathllap{\begin{aligned}    \de\yF_{\yc}/\de\yp&=0\text{:} \\
%         \de\yF_{\yc}/\de\ym&=0\text{:}\\ \de\yF_{\yc}/\de\yl&=0\text{:}\end{aligned}}\qquad}\empheqlbrace]{align}
%     \label{eq:con_p}
% %    \de\yF_{\yc}/\de\yp &\equiv
%     -\ln\yp + \ln\yq + \yl\yM + \ym\yu &=0,\\
%     \label{eq:con_u}
% %    \de\yF_{\yc}/\de\ym &\equiv
%     \yu\yp-1 &=0,\\
%     \label{eq:con_l}
%     %\de\yF_{\yc}/\de\yl &\equiv
%     \yM\yp-\yc &=0.
%   \end{empheq}
%%%%
% \begin{empheq}[box=\widefbox]{equation}
%   \label{eq:maxent_question}
%   \p\bigl[\yE{N+1}{k} \bigcond \tsum\yo\yf{N}\in\yA, \yM\bigr] = \mathord{?}
% \end{empheq}



% \[
%   \begin{tikzcd}
%       M_{n,n}(\CC) \arrow{r}{R'_{a}(\Hat{U})} & M_{n,n}(\CC)
%     \\
%     L(\mathcal{H}) \arrow{r}{\Hat{U}} \arrow[swap]{d}{R_*}\arrow[swap]{u}{R'_*} & L(\mathcal{H}) \arrow{d}{R_*}\arrow{u}{R'_*} \\
%       M_{n,n}(\CC) \arrow{r}{R_{a}(\Hat{U})} & M_{n,n}(\CC)
%   \end{tikzcd}
% \]

% \[
%   \begin{tikzcd}
%       \CC^n \arrow{r}{R'_*(A)} & \CC^n
%     \\
%     \mathcal{H} \arrow{r}{A} \arrow[swap]{d}{R}\arrow[swap]{u}{R'} & \mathcal{H} \arrow{d}{R}\arrow{u}{R'} \\
%       \CC^n \arrow{r}{R_*(A)} & \CC^n
%   \end{tikzcd}
% \]


% \[
%   \begin{tikzcd}
%     \mathcal{H} \arrow{r}{A} \arrow[swap]{d}{R} & \mathcal{H} \arrow{d}{R} \\
%       \CC^n \arrow{r}{R_*(A)} & \CC^n
%   \end{tikzcd}
% \]

%%\setlength{\intextsep}{0ex}% with wrapfigure
%%\setlength{\columnsep}{0ex}% with wrapfigure
%\begin{figure}[p!]% with figure
%\begin{wrapfigure}{r}{0.4\linewidth} % with wrapfigure
%  \centering\includegraphics[trim={12ex 0 18ex 0},clip,width=\linewidth]{maxent_saddle.png}\\
%\caption{caption}\label{fig:comparison_a5}
%\end{figure}% exp_family_maxent.nb


%%%%%%%%%%%%%%%%%%%%%%%%%%%%%%%%%%%%%%%%%%%%%%%%%%%%%%%%%%%%%%%%%%%%%%%%%%%%
%%% Acknowledgements
%%%%%%%%%%%%%%%%%%%%%%%%%%%%%%%%%%%%%%%%%%%%%%%%%%%%%%%%%%%%%%%%%%%%%%%%%%%% 
\iffalse
\begin{acknowledgements}
  \ldots to Mari \amp\ Miri for continuous encouragement and affection, and
  to Buster Keaton and Saitama for filling life with awe and inspiration.
  To the developers and maintainers of \LaTeX, Emacs, AUC\TeX, Open Science
  Framework, R, Python, Inkscape, Sci-Hub for making a free and impartial
  scientific exchange possible.
  % Our work was supported by the Trond Mohn Research Foundation, grant number BFS2018TMT07
%\rotatebox{15}{P}\rotatebox{5}{I}\rotatebox{-10}{P}\rotatebox{10}{\reflectbox{P}}\rotatebox{-5}{O}.
%\sourceatright{\autanet}
\mbox{}\hfill\autanet
\end{acknowledgements}
\fi

%%%%%%%%%%%%%%%%%%%%%%%%%%%%%%%%%%%%%%%%%%%%%%%%%%%%%%%%%%%%%%%%%%%%%%%%%%%%
%%% Appendices
%%%%%%%%%%%%%%%%%%%%%%%%%%%%%%%%%%%%%%%%%%%%%%%%%%%%%%%%%%%%%%%%%%%%%%%%%%%% 
%\clearpage
\bigskip
% %\renewcommand*{\appendixpagename}{Appendix}
% %\renewcommand*{\appendixname}{Appendix}
% %\appendixpage
% \appendix

%%%%%%%%%%%%%%%%%%%%%%%%%%%%%%%%%%%%%%%%%%%%%%%%%%%%%%%%%%%%%%%%%%%%%%%%%%%%
%%% Bibliography
%%%%%%%%%%%%%%%%%%%%%%%%%%%%%%%%%%%%%%%%%%%%%%%%%%%%%%%%%%%%%%%%%%%%%%%%%%%% 
\renewcommand*{\finalnamedelim}{\addcomma\space}
\defbibnote{prenote}{{\footnotesize (\enquote{de $X$} is listed under D,
    \enquote{van $X$} under V, and so on, regardless of national
    conventions.)\par}}
% \defbibnote{postnote}{\par\medskip\noindent{\footnotesize% Note:
%     \arxivp \mparcp \philscip \biorxivp}}

\printbibliography[prenote=prenote%,postnote=postnote
]

\end{document}

%%%%%%%%%%%%%%%%%%%%%%%%%%%%%%%%%%%%%%%%%%%%%%%%%%%%%%%%%%%%%%%%%%%%%%%%%%%%
%%% Cut text (won't be compiled)
%%%%%%%%%%%%%%%%%%%%%%%%%%%%%%%%%%%%%%%%%%%%%%%%%%%%%%%%%%%%%%%%%%%%%%%%%%%% 


%%% Local Variables: 
%%% mode: LaTeX
%%% TeX-PDF-mode: t
%%% TeX-master: t
%%% End: 
